\section*{الفصل \ref{ch:g12:diffeq} --- المعادلات التفاضلية}

\begin{solution}{exo:g12:diffeq:1}
\emph{(أ)} $y = C\eu^{3x}$؛ و $y(0) = 2$ يعطي $y = 2\eu^{3x}$.

\emph{(ب)} $y' = -\frac12 y$، إذن $y = C\eu^{-x/2}$؛ و $y(0) = -1$ يعطي
$y = -\eu^{-x/2}$.

\emph{(ج)} التوازن $y_p = 5$؛ و $y = C\eu^{-x} + 5$؛ و $y(0) = 0$ يعطي
$C = -5$: $y = 5\left(1 - \eu^{-x}\right)$.
\end{solution}

\begin{solution}{exo:g12:diffeq:2}
$N' = kN$، إذن $N(t) = N_0 \eu^{kt}$. والتضاعف في $3$ ساعات يعني
$\eu^{3k} = 2$، \emph{أي}
\[
k = \frac{\ln 2}{3} \approx 0.231\ \text{h}^{-1}.
\]
\end{solution}

\begin{solution}{exo:g12:diffeq:3}
$y_p'(x) = \eu^x + x\eu^x = y_p(x) + \eu^x$: إذن $y_p$ حل خاص
(وهذه هي حالة الرنين في \cref{met:g12:diffeq:particular}). وحسب
\cref{prop:g12:diffeq:general}، تكون الحلول
$y = C\eu^{x} + x\eu^{x} = (C + x)\,\eu^x$، حيث $C \in \R$.
\end{solution}

\begin{solution}{exo:g12:diffeq:4}
جرّب $y_p = \beta\eu^x$: فإن $\beta\eu^x = -2\beta\eu^x + \eu^x$ يعطي
$3\beta = 1$، إذن $y_p = \frac13\eu^x$. والحل العام
$y = C\eu^{-2x} + \frac13\eu^{x}$؛ والشرط $y(0) = 1$ يعطي
$C = \frac23$:
\[
y(x) = \frac{2}{3}\,\eu^{-2x} + \frac{1}{3}\,\eu^{x}.
\]
\end{solution}

\begin{solution}{exo:g12:diffeq:5}
$N(t) = N_0\,\eu^{-\lambda t}$ مع $\lambda = \frac{\ln 2}{5730}$
(\cref{ex:g12:exp:halflife}). ونحل $\eu^{-\lambda t} = 0.6$:
\[
t = \frac{\ln(1/0.6)}{\lambda} = 5730\,\frac{\ln(5/3)}{\ln 2}
\approx 5730 \times \frac{0.5108}{0.6931} \approx 4220 \text{ سنة}.
\]
\end{solution}

\begin{solution}{exo:g12:diffeq:6}
\emph{1.} يدخل الملح بمعدل $0.2 \times 5 = 1$ kg/min. ويحمل التدفق الخارج
تركيزًا $\frac{m}{100}$ kg/L بمعدل $5$ L/min، \emph{أي} $\frac{m}{20}$
kg/min. ومنه $m' = 1 - \frac{m}{20}$.

\emph{2.} التوازن $m_p = 20$؛ و $m(t) = 20 + C\eu^{-t/20}$، و
$m(0) = 0$ يعطي $C = -20$:
\[
m(t) = 20\left(1 - \eu^{-t/20}\right) \xrightarrow[t\to+\infty]{} 20 \text{ kg}.
\]
وعلى \omterm{def:g10:stats:quartiles}{المدى} الطويل يساوي تركيز الخزان تركيز المحلول الداخل:
$0.2$ kg/L $\times$ $100$ L $= 20$ kg.
\end{solution}

\begin{solution}{exo:g12:diffeq:7}
\emph{1.} $v' = g - \frac{k}{m} v$ مع $\frac km = \frac{16}{80} = 0.2$.
والتوازن $v_\infty = \frac{mg}{k} = \frac{9.8}{0.2} = 49$ m/s؛
و $v(t) = 49 + C\eu^{-0.2t}$، و $v(0) = 0$ يعطي
\[
v(t) = 49\left(1 - \eu^{-0.2 t}\right).
\]

\emph{2.} السرعة الحدية $49$ m/s ($\approx 176$ km/h). ونريد
$1 - \eu^{-0.2t} = 0.95$، \emph{أي} $\eu^{-0.2t} = 0.05$:
\[
t = \frac{\ln 20}{0.2} \approx \frac{3.00}{0.2} \approx 15 \text{ s}.
\]
\end{solution}

\begin{solution}{exo:g12:diffeq:8}
\emph{1.} $z = \frac1y$ يعطي $z' = -\frac{y'}{y^2}
= -\frac{y(1-y)}{y^2} = -\frac{1-y}{y} = -\frac1y + 1 = -z + 1$.

\emph{2.} التوازن $z_p = 1$، إذن $z(t) = 1 + C\eu^{-t}$. ومن
$y(0) = \frac{1}{10}$، يكون $z(0) = 10$، إذن $C = 9$ و
\[
y(t) = \frac{1}{1 + 9\,\eu^{-t}} .
\]

\emph{3.} لما $t \to +\infty$، يكون $9\eu^{-t} \to 0$ و $y(t) \to 1$. \omterm{def:g10:functions:graph}{والمنحنى}
هو \omterm{def:g10:functions:graph}{المنحنى} اللوجستي الكلاسيكي على شكل S (\emph{السيني}): نمو بطيء
في البداية ($y$ صغير و $y' \approx y$)، وأسرع نمو عند $y = \frac12$
(حيث يكون $y' = y(1-y)$ أعظميًا)، ثم إشباع نحو سعة
\omterm{def:g10:proba:distribution}{الاحتمال} $1$.
\end{solution}

\begin{solution}{exo:g12:diffeq:9}
حسب \cref{thm:g12:diffeq:affine}، $y(x) = C\eu^{ax} - \frac ba$. ومنه
\[
u_{n+1} = C\eu^{a(n+1)} - \frac ba
= \eu^{a}\left(C\eu^{an} - \frac ba\right) + \frac ba\left(\eu^a - 1\right)
= q\,u_n + r
\]
مع $q = \eu^{a}$ و $r = \frac{b}{a}\left(\eu^{a} - 1\right)$.
والشرط $\abs q < 1$ يعني $\eu^a < 1$، \emph{أي} $a < 0$: وهو بالضبط
الشرط الذي \omterm{def:g12:seq:limit}{تتقارب} تحته حلول \omterm{def:g12:diffeq:ode}{المعادلة التفاضلية}
إلى التوازن $-\frac ba$ --- وفعلًا فإن \omterm{pb:g10:functions:1}{النقطة الثابتة}
للعلاقة التراجعية هي $\frac{r}{1 - q} = -\frac{b}{a}$.
\end{solution}

\begin{solution}{pb:g12:diffeq:1}
\textbf{1.} $y = 2\eu^{3t}$. ومن أجل الثانية: التوازن $3$،
إذن $y = 3 + C\eu^{-2t}$ مع $y(0) = 0$: أي $C = -3$:
$y = 3\left(1 - \eu^{-2t}\right)$.

\textbf{2.} $\left(y\eu^{-at}\right)' = y'\eu^{-at} -
ay\eu^{-at} = (y' - ay)\eu^{-at} = 0$: فالجداء مقدار
ثابت $C$، إذن $y = C\eu^{at}$ --- فلا حل يفلت.

\textbf{3.} التوازن: $y \equiv 3$. وكل حل آخر هو
$3 + C\eu^{-2t}$ حيث $C \neq 0$: فتموت \omterm{thm:g12:exp:existence}{الدالة الأسية} وينزلق
الحل إلى $3$ --- وهي \omterm{pb:g10:functions:1}{نقطة ثابتة} مستقرة، التوأم
المتصل للنقط الثابتة التراجعية في
\cref{pb:g11:seq:1}.

\textbf{4.} $y = y_0\eu^{-t/\tau}$ ينتصف عندما يكون
$\eu^{-t/\tau} = \frac12$: أي $t = \tau\ln 2$.

\textbf{5.} كل الحلول انسحابات شاقولية للتفكك
نحو $3$: فالتي تبدأ فوقه تهبط إلى $3$، والتي تبدأ تحته
ترتفع إلى $3$، والحل الثابت $3$ يقف ساكنًا --- إنه
قِمع حول التوازن.

\textbf{6.} $\lambda = \frac{\ln 2}{5730} \approx 1.21 \times
10^{-4}$ في السنة.

\textbf{7.} $\eu^{-\lambda t} = 0.2$:
$t = \frac{\ln 5}{\lambda} \approx 13\,300$ سنة.

\textbf{8.} $t = \frac{\ln(1/0.15)}{\lambda} \approx 15\,700$
سنة: فثيران لاسكو من أواخر العصر الجليدي.

\textbf{9.} $t = \frac{\ln(1/0.53)}{\lambda} \approx 5\,250$
سنة: أي ضمن عمر إنسان من قيمة علماء الآثار ---
فالساعة تعمل.

\textbf{10.} عشرة أنصاف أعمار تترك $2^{-10} \approx 0.1\,\%$:
أي على حافة القياس. وبعد $65$ مليون سنة يكون
الجزء $2^{-65\,000\,000/5\,730} \approx 2^{-11\,344}$:
فلا تبقى ذرة من المخزون الأصلي في أي أحفورة ---
وتُؤرَّخ الديناصورات بساعات أبطأ (البوتاسيوم--الأرغون
وأمثاله).

\textbf{11.} $52\,\%$ تعطي $\approx 5\,405$ سنة،
و $54\,\%$ تعطي $\approx 5\,095$: أي إن $\pm 1\,\%$ من الكيمياء
تصير نحو $\pm 155$ سنة من التاريخ --- فالدقة في
المخبر دقة في لوحة المتحف.

\textbf{12.} $z = T - T_a$ تحقق $z' = -kz$:
$z = (T_0 - T_a)\eu^{-kt}$، و $T = T_a + z$.

\textbf{13.} الفجوة الابتدائية $T_0 - T_a = 70$؛ وبعد خمس
دقائق تكون الفجوة $70 - 20 = 50$، إذن $50 = 70\eu^{-5k}$:
أي $k = \frac{\ln(70/50)}{5} \approx 0.067$ في الدقيقة. وتصير قابلة للشرب عند:
$55 - 20 = 35 = 70\eu^{-kt}$:
$t = \frac{\ln 2}{k} \approx 10.3$ دقيقة.

\textbf{14.} معدل الفقد يتناسب مع الفجوة
$T - T_a$: فالقهوة الحارقة تنزف الحرارة أسرع. وإضافة
الحليب \emph{الآن} تقلّص الفجوة فورًا، فيُفقد حرّ أقل
خلال العشر دقائق؛ أما إضافته في النهاية فتترك
القهوة تبرد بأقصى سرعة أولًا. فمن أجل أدفأ فنجان: الحليب
أولًا. (والمضيفون العجولون يقلبون الترموديناميك.)

\textbf{15.} الفجوات عن المحيط: عند \omterm{prop:g10:coordgeom:midpoint}{منتصف} الليل $10$، وبعد ساعة
$8$: إذن $\eu^{-k} = 0.8$:
$k = \ln\frac{10}{8} \approx 0.223$ في الساعة. وعند الوفاة كانت الفجوة
$17$؛ وقد تفككت إلى $10$ عند \omterm{prop:g10:coordgeom:midpoint}{منتصف} الليل:
$s = \frac{\ln(17/10)}{k} \approx 2.4$ ساعة. فساعة الوفاة:
نحو $21{:}40$.

\textbf{16.} غرفة مدفَّأة أو ذات تيار هواء ($T_a$ غير ثابتة)
تحرف الساعة في أي من الجهتين؛ والملابس أو كتلة الجسم تغيّر $k$
(وجداول الطبيب الشرعي تصحّح لذلك) --- و $k$ الخاطئ يضرب
\omterm{def:g10:numbers:interval}{الفترة} كلها بمعامل؛ وجثة نُقلت من مكان آخر تعيد ضبط
$T_a$ في \omterm{prop:g10:coordgeom:midpoint}{منتصف} التفكك، فتزيّف وفاة أبكر أو أمتأخر. \omterm{def:g10:algebra:equation}{فالمعادلة}
أمينة؛ والفرضيات هي التي تحمل الخطر.

\textbf{17.} $z = \frac1y$ تطيع $z' = 1 - z$:
$z = 1 + 9\eu^{-t}$ (من $z(0) = 10$)، إذن
$y = \frac{1}{1 + 9\eu^{-t}}$. والانعطاف عند $y = \frac12$:
$9\eu^{-t} = 1$: $t = \ln 9 \approx 2.2$ --- وهو تاريخ انقلاب
الوباء، مباشرة من الصيغة.

\textbf{18.} السرعة الحدية: $v' = 0$ عند $v = 20$ m/s.
والحل: $v(t) = 20\left(1 - \eu^{-t/2}\right)$. ثم
$0.95$: $\eu^{-t/2} = 0.05$: $t = 2\ln 20 \approx 6$ s ---
فقطرات المطر تبلغ سرعتها المستقرة في ثوانٍ، ولهذا
لا يقتل المطر.

\textbf{19.} الحالة المستقرة: $c = 8$. ونصفها:
$4 = 8\left(1 - \eu^{-t/2}\right)$: $\eu^{-t/2} = \frac12$:
$t = 2\ln 2 \approx 1.4$ ساعة. فالمحلول الوريدي هو التوأم المتصل
للقرض: دخل ثابت، وخروج تناسبي ---
و \cref{exo:g12:diffeq:9} يجعل القاموس مضبوطًا.

\textbf{20.} التفكك ($a < 0$ و $b = 0$)، والتبرّد ($y = T - T_a$)،
والسقوط ($a < 0$ و $b > 0$: اقتراب من التوازن من
تحت)، والجرعات (المثل نفسه، في وريد): \omterm{def:g10:algebra:equation}{معادلة} واحدة، وأربعة عوالم.
والحلقة: صُغ قانون التغيّر؛ واكتب \omterm{def:g10:algebra:equation}{المعادلة}؛ وحلّ
(\cref{met:g12:diffeq:solve})؛ وعاير الثوابت على
القياسات؛ وتنبّأ؛ ثم دقّق الفرضيات (السؤال
16). وأما التذبذب فيحتاج إلى $y'' = -\omega^2 y$ --- وقد لُقي وحُلّ و
أُطلق متأرجحًا في \cref{pb:g12:trigo:1}.
\end{solution}
